\section{Details of Considered Processes}
\label{app:details of considered processes}
    SEDD~\citep{lou2024discrete} introduced a general framework for discrete diffusion models based on an arbitrary diffusion matrix $Q_t$, allowing the modeling of a broad class of forward processes. 
    Subsequent works, including MDLM~\citep{sahoo2024simple}, UDLM~\citep{schiff2024discreteguidance}, and Duo~\citep{sahoo2025the}, build upon this framework by selecting specific diffusion dynamics and adopting distinct parameterizations and training objectives.
    In most settings, the diffusion matrix is expressed in the form $Q_t = \sigma_t Q$, where $Q$ is a fixed base matrix defining the structure of the process, and $\sigma_t$ is a time-dependent scalar that governs the diffusion rate. Defining the cumulative diffusion schedule as $\overline{\sigma}_t \coloneqq \int_{0}^t \sigma_s \, ds$, one can further express the cumulative transition dynamics using the matrix exponential $\exp(\overline{\sigma}_t Q)$. Two principal classes of diffusion processes have emerged in practice: the \textit{absorbing} and \textit{uniform} processes, instantiated by selecting $Q = Q_{\text{abs}}$ and $Q = Q_{\text{uni}}$, respectively. The purpose of this section is to formally introduce and characterize these two diffusion processes.
    \subsection{Absorbing process}
        \label{app:absorbing process}
        Without loss of generality, we define the mask token as $m \in \mathcal{V}$ such that $m_N = 1$. An important special case of discrete diffusion is the \textit{absorbing process}, in which the terminal distribution converges to a delta distribution concentrated on a designated mask token: $$\pi = m.$$ This process is realized by choosing a specific base diffusion matrix~\citep{lou2024discrete}:
        \begin{equation}
            Q_{\text{abs}} =
            \begin{bmatrix}
                -1 & 0 & \cdots & 0 & 0 \\
                0 & -1 & \cdots & 0 & 0 \\
                \vdots & \vdots & \ddots & \vdots & \vdots \\
                0 & 0 & \cdots & -1 & 0 \\
                1 & 1 & \cdots & 1 & 0
            \end{bmatrix}.
        \end{equation}
        This choice induces the time-dependent diffusion matrix
        \begin{equation}
            \label{eq: absorbing diffusion matrix}
            Q_t = \sigma_t Q_{\text{abs}} =
            \begin{bmatrix}
                -\sigma_t & 0 & \cdots & 0 & 0 \\
                0 & -\sigma_t & \cdots & 0 & 0 \\
                \vdots & \vdots & \ddots & \vdots & \vdots \\
                0 & 0 & \cdots & -\sigma_t & 0 \\
                \sigma_t & \sigma_t & \cdots & \sigma_t & 0
            \end{bmatrix}.
        \end{equation}
        
        The corresponding cumulative transition matrix is given by the matrix exponential
        \begin{equation}
            \exp(\overline{\sigma}_t Q_{\text{abs}}) =
            \begin{bmatrix}
                \exp(-\overline{\sigma}_t) & 0 & \cdots & 0 & 0 \\
                0 & \exp(-\overline{\sigma}_t) & \cdots & 0 & 0 \\
                \vdots & \vdots & \ddots & \vdots & \vdots \\
                0 & 0 & \cdots & \exp(-\overline{\sigma}_t) & 0 \\
                1 - e^{-\overline{\sigma}_t} & 1 - e^{-\overline{\sigma}_t} & \cdots & 1 - e^{-\overline{\sigma}_t} & 1
            \end{bmatrix}.
        \end{equation}
        
        In MDLM~\citep[Appendix C]{sahoo2024simple}, an alternative parameterization of the diffusion schedule is employed by defining $\alpha_t \coloneqq \exp(-\overline{\sigma}_t)$. Under this notation, the diffusion matrix can be expressed as
        \begin{equation}
            Q_t = -\frac{\alpha_t'}{\alpha_t} Q_{\text{abs}} =
            \begin{bmatrix}
                \frac{\alpha_t'}{\alpha_t} & 0 & \cdots & 0 & 0 \\
                0 & \frac{\alpha_t'}{\alpha_t} & \cdots & 0 & 0 \\
                \vdots & \vdots & \ddots & \vdots & \vdots \\
                0 & 0 & \cdots & \frac{\alpha_t'}{\alpha_t} & 0 \\
                -\frac{\alpha_t'}{\alpha_t} & -\frac{\alpha_t'}{\alpha_t} & \cdots & -\frac{\alpha_t'}{\alpha_t} & 0
            \end{bmatrix},
        \end{equation}
        and in the same way the cumulative diffusion matrix can be equivalently expressed as
        \begin{equation}
            \exp(\overline{\sigma}_t Q_{\text{abs}}) =
            \begin{bmatrix}
                \alpha_t & 0 & \cdots & 0 & 0 \\
                0 & \alpha_t & \cdots & 0 & 0 \\
                \vdots & \vdots & \ddots & \vdots & \vdots \\
                0 & 0 & \cdots & \alpha_t & 0 \\
                1 - \alpha_t & 1 - \alpha_t & \cdots & 1 - \alpha_t & 1
            \end{bmatrix}.
        \end{equation}
    \subsection{Uniform process}
        \label{app:uniform process}
        Another important special case of discrete diffusion is the \textit{uniform process}, in which the terminal distribution is uniform: 
        $$\pi = \frac{1}{N} \vec{1}.$$
        This process is realized by selecting the following base diffusion matrix~\citep{lou2024discrete}:
        \begin{equation}
            Q_{\text{uni}} =
            \begin{bmatrix}
                1 - N & 1 & \cdots & 1 \\
                1 & 1 - N & \cdots & 1 \\
                \vdots & \vdots & \ddots & \vdots \\
                1 & 1 & \cdots & 1 - N
            \end{bmatrix}.
        \end{equation}
        Given a diffusion rate $\sigma_t$, this leads to the time-dependent diffusion matrix:
        \begin{equation}
        \label{eq: uniform diffusion matrix}
            Q_t = \sigma_t Q_{\text{uni}} =
            \begin{bmatrix}
                \sigma_t(1 - N) & \sigma_t & \cdots & \sigma_t \\
                \sigma_t & \sigma_t(1 - N) & \cdots & \sigma_t \\
                \vdots & \vdots & \ddots & \vdots \\
                \sigma_t & \sigma_t & \cdots & \sigma_t(1 - N)
            \end{bmatrix}.
        \end{equation}
        The corresponding cumulative transition matrix is given by the matrix exponential:
        \begin{equation}
            \exp(\overline{\sigma}_t Q_{\text{uni}}) =
            \begin{bmatrix}
                e^{-\overline{\sigma}_t N} + \frac{1 - e^{-\overline{\sigma}_t N}}{N} & \frac{1 - e^{-\overline{\sigma}_t N}}{N} & \cdots & \frac{1 - e^{-\overline{\sigma}_t N}}{N} \\
                \frac{1 - e^{-\overline{\sigma}_t N}}{N} & e^{-\overline{\sigma}_t N} + \frac{1 - e^{-\overline{\sigma}_t N}}{N} & \cdots & \frac{1 - e^{-\overline{\sigma}_t N}}{N} \\
                \vdots & \vdots & \ddots & \vdots \\
                \frac{1 - e^{-\overline{\sigma}_t N}}{N} & \frac{1 - e^{-\overline{\sigma}_t N}}{N} & \cdots & e^{-\overline{\sigma}_t N} + \frac{1 - e^{-\overline{\sigma}_t N}}{N}
            \end{bmatrix}.
        \end{equation}
        
        In UDLM~\citep{schiff2024discreteguidance}, an alternative parameterization of the diffusion schedule is adopted by defining $\alpha_t \coloneqq \exp(-\overline{\sigma}_t N)$. Under this parameterization, the diffusion matrix becomes:
        \begin{equation}
            Q_t = -\frac{\alpha_t'}{N\alpha_t} Q_{\text{uni}} =
            \begin{bmatrix}
                - \frac{\alpha_t'}{N\alpha_t}(1 - N) & -\frac{\alpha_t'}{N\alpha_t} & \cdots & -\frac{\alpha_t'}{N\alpha_t} \\
                -\frac{\alpha_t'}{N\alpha_t} & -\frac{\alpha_t'}{N\alpha_t}(1 - N) & \cdots & -\frac{\alpha_t'}{N\alpha_t} \\
                \vdots & \vdots & \ddots & \vdots \\
                -\frac{\alpha_t'}{N\alpha_t} & -\frac{\alpha_t'}{N\alpha_t} & \cdots & -\frac{\alpha_t'}{N\alpha_t}(1 - N)
            \end{bmatrix}.
        \end{equation}
        
        Likewise, defining $\beta_t = \frac{1 - \alpha_t}{N}$, the cumulative transition matrix under this parameterization is expressed as:
        \begin{equation}
            \exp(\overline{\sigma}_t Q_{\text{uni}}) =
            \begin{bmatrix}
                \alpha_t + \beta_t & \beta_t & \cdots & \beta_t \\
                \beta_t & \alpha_t + \beta_t & \cdots & \beta_t \\
                \vdots & \vdots & \ddots & \vdots \\
                \beta_t & \beta_t & \cdots & \alpha_t + \beta_t
            \end{bmatrix}.
        \end{equation}

    \subsection{Backward Diffusion Details}
    \label{app:backward_diffusion}
    If the backward diffusion matrix $\overline{Q}_t$ is known, the reverse process can be simulated using numerical methods such as Euler's~\citep{lou2024discrete}. It can be derived from the forward diffusion matrix $Q_t$ using the following relationship~\citep{kelly2011reversibility, campbell2022continuous, sun2022score}:
    \begin{equation*}
      \langle y, x\rangle_{\overline{Q}_{t}} = \frac{p_t(y)}{p_t(x)} \langle x, y \rangle_{Q_t}, \quad 
      \langle x, x\rangle_{\overline{Q}_{t}} = - \sum_{y \neq x} \langle y, x \rangle_{\overline{Q}_{t}},
    \end{equation*}
    where $y, x \in \VC$. The main challenge is to compute the ratio $\frac{p_t(y)}{p_t(x)}$, which depends on unknown intermediate distributions and cannot be computed directly.
